
\clearpage
\chapter{עבודות קשורות: מודלי רצפים}

\section{מודלי מרחב מצב (\textenglish{SSM})}

\textenglish{S4} \cite{gu2021s4} הציג מודל מרחב מצב (\textenglish{Structured State Space}) המאחד גישות \textenglish{RNN} וקונבולוציה: בזמן אימון, \textenglish{S4} מחושב כקונבולוציה גלובלית יעילה; בזמן אינפרנס, כ-\textenglish{RNN} עם מצב מצטבר. \textenglish{S4} הציג ביצועים מרשימים על \textenglish{Long Range Arena} אך נחות מ-\textenglish{xLSTM} על חיזוי טורי זמן קצר-ארוך \cite{beck2024xlstm}.

\section{מודלים ליניאריים}

\textenglish{DLinear} \cite{zeng2023dlinear} הפתיע את הקהילה: \textenglish{1-layer MLP} עם פירוק מגמה-עונתיות הדגים ביצועים תחרותיים לעומת \textenglish{Transformer} על מספר \textenglish{benchmarks}. הממצא הזה הוביל לשאלה מחקרית מרכזית: האם \textenglish{Transformer} באמת מתאים לחיזוי טורי זמן? \textenglish{xLSTM} עונה על שאלה זו בחיוב, אך מספק מנגנוני זיכרון חזקים יותר מ-\textenglish{DLinear}.

\section{ארכיטקטורות מבוססות \textenglish{Attention} ספציפיות לזמן}

\textenglish{Pyraformer} \cite{liu2022pyraformer} מציג מבנה פירמידאלי להפחתת מורכבות הקשב מ-\textenglish{O(n²)} ל-\textenglish{O(n)}. \textenglish{FEDformer} \cite{zhou2022fedformer} עובד בתחום תדר לניצול תכונות ספקטרליות. שתי הגישות יעילות בחיסכון עלות חישובית אך משלמות מחיר בדיוק.

\section{השוואה רחבה}

הספרות מראה ש-\textenglish{xLSTM} מוביל ב: (א) חיזוי לטווח בינוני (\textenglish{H=96-336}); (ב) מערכי נתונים עם תלויות זמניות חזקות (למשל, צריכת חשמל). \textenglish{PatchTST} \cite{nie2022patchtst} קרוב מאד ועדיף ב-\textenglish{few-shot} עקב \textenglish{ICL}. לחיזוי ארוך (\textenglish{H=720}) על מערכי נתונים סטוכסטיים, הפרשים קטנים ואינם בהכרח מובהקים סטטיסטית.
